
\documentclass{article}
\usepackage{colortbl}
\usepackage{makecell}
\usepackage{multirow}
\usepackage{supertabular}

\begin{document}

\newcounter{utterance}

\centering \large Interaction Transcript for game `hot\_air\_balloon', experiment `air\_balloon\_survival\_sr\_complexity\_hard', episode 0 with gemma{-}4{-}26b{-}a4b{-}it.
\vspace{24pt}

{ \footnotesize  \setcounter{utterance}{1}
\setlength{\tabcolsep}{0pt}
\begin{supertabular}{c@{$\;$}|p{.15\linewidth}@{}p{.15\linewidth}p{.15\linewidth}p{.15\linewidth}p{.15\linewidth}p{.15\linewidth}}
    \# & $\;$A & \multicolumn{4}{c}{Game Master} & $\;\:$B\\
    \hline

    \theutterance \stepcounter{utterance}  
    & & \multicolumn{4}{p{0.6\linewidth}}{
        \cellcolor[rgb]{0.9,0.9,0.9}{
            \makecell[{{p{\linewidth}}}]{
                \texttt{\tiny{[P1$\langle$GM]}}
                \texttt{Učestvujete u kolaborativnoj igri pregovaranja.} \\
\\ 
\texttt{Zajedno sa drugim učesnikom, morate da se dogovorite o jednom jedinom skupu predmeta koji će biti zadržan. Svako od vas ima sopstveni pogled na to koliko vam je svaki predmet važan (važnost). Ne znate kako drugi učesnik vrednuje predmete. Pored toga, dat vam je napor koji svaki predmet zahteva.} \\
\texttt{Možete se dogovoriti o skupu samo ako ukupan napor izabranih predmeta ne prelazi zajedničko ograničenje:} \\
\\ 
\texttt{LIMIT = 4734} \\
\\ 
\texttt{Evo pojedinačnih vrednosti napora po predmetu:} \\
\\ 
\texttt{Item effort = \{"C08": 192, "C55": 325, "A43": 812, "A67": 868, "B42": 593, "A62": 915, "B19": 944, "C43": 311, "C73": 252, "C10": 343, "B36": 104, "A18": 558, "C61": 627, "C67": 593, "C42": 827, "C70": 611, "B40": 95, "C39": 251, "A21": 226, "C16": 21\}} \\
\\ 
\texttt{Evo vašeg ličnog pogleda na važnost svakog predmeta:} \\
\\ 
\texttt{Item importance values = \{"C08": 286, "C55": 419, "A43": 906, "A67": 962, "B42": 687, "A62": 1009, "B19": 1038, "C43": 405, "C73": 346, "C10": 437, "B36": 198, "A18": 652, "C61": 721, "C67": 687, "C42": 921, "C70": 705, "B40": 189, "C39": 345, "A21": 320, "C16": 115\}} \\
\\ 
\texttt{Cilj:} \\
\\ 
\texttt{Vaš cilj je da pregovorima postignete zajednički skup predmeta koji vam koristi što je više moguće (tj. maksimizuje ukupnu važnost za VAS), uz poštovanje LIMIT{-}a. Niste obavezni da u svakoj poruci date PROPOSAL {-} možete i samo da pregovarate. Sve taktike su dozvoljene!} \\
\\ 
\texttt{Protokol interakcije:} \\
\\ 
\texttt{U poruci smete da koristite samo sledeće strukturirane formate:} \\
\\ 
\texttt{PREDLOG: \{'A', 'B', 'C', …\}} \\
\texttt{Predložite da se zadrže tačno ti predmeti.} \\
\texttt{ODBIJAM: \{'A', 'B', 'C', …\}} \\
\texttt{Izričito odbijte protivnikov predlog.} \\
\texttt{ARGUMENT: \{'...'\}} \\
\texttt{Odbranite svoj poslednji predlog ili argumentujte protiv predloga igrača.} \\
\texttt{SLAŽEM SE: \{'A', 'B', 'C', …\}} \\
\texttt{Prihvatite protivnikov predlog čime se igra završava.} \\
\\ 
\\ 
\texttt{Pravila:} \\
\\ 
\texttt{Možete dati AGREE samo na predlog koji je druga strana evidentirala putem PROPOSAL.} \\
\texttt{Možete dati REFUSE samo na predlog koji je druga strana evidentirala putem PROPOSAL.} \\
\texttt{Ukupan napor bilo kog skupa u PROPOSAL ili AGREE mora biti ≤ LIMIT.} \\
\texttt{NE otkrivajte svoje skrivene ocene važnosti.} \\
\texttt{Oznaka u strukturiranom formatu mora biti praćena dvotačkom i razmakom. Argument mora biti python skup koji sadrži 0 ili više stringova.} \\
\texttt{Dakle, mora biti u obliku TAG: \{...\}} \\
\texttt{Strogo se pridržavajte protokola interakcije i NE pišite ništa izvan date strukture.} \\
\texttt{Igra se završava kada jedna strana da AGREE na PROPOSAL koji je dao drugi igrač.} \\
\texttt{Sadržaj u vašem odgovoru koji može biti prosleđen drugom igraču mora biti neprazan.} \\
\texttt{Aktivni su samo predlozi koji su evidentirani putem strukture formata PROPOSAL i koji nisu odbijeni putem REFUSE.} \\
\\ 
            }
        }
    }
    & \\ \\

    \theutterance \stepcounter{utterance}  
    & \multicolumn{4}{p{0.6\linewidth}}{
        \cellcolor[rgb]{0.9,0.9,0.9}{
            \makecell[{{p{\linewidth}}}]{
                \texttt{\tiny{[P1$\rangle$GM]}}
                \texttt{PREDLOG: \{'A62', 'B19', 'A67', 'A43', 'C42', 'C61', 'C70', 'B42', 'C67', 'A18'\}} \\
            }
        }
    }
    & & \\ \\

    \theutterance \stepcounter{utterance}  
    & & \multicolumn{4}{p{0.6\linewidth}}{
        \cellcolor[rgb]{0.9,0.9,0.9}{
            \makecell[{{p{\linewidth}}}]{
                \texttt{\tiny{[GM$\rangle$P2]}}
                \texttt{Učestvujete u kolaborativnoj igri pregovaranja.} \\
\\ 
\texttt{Zajedno sa drugim učesnikom, morate se dogovoriti o jednom jedinstvenom skupu predmeta koji će biti zadržan. Svako od vas ima sopstveni pogled na to koliko vam je svaki predmet važan (važnost). Ne znate kako drugi učesnik vrednuje predmete. Pored toga, dobijate i napor koji svaki predmet zahteva.} \\
\texttt{Možete se dogovoriti o skupu samo ako ukupan napor izabranih predmeta ne prelazi zajedničko ograničenje:} \\
\\ 
\texttt{LIMIT = 4734} \\
\\ 
\texttt{Evo pojedinačnih vrednosti napora po predmetu:} \\
\\ 
\texttt{Item effort = \{"C08": 192, "C55": 325, "A43": 812, "A67": 868, "B42": 593, "A62": 915, "B19": 944, "C43": 311, "C73": 252, "C10": 343, "B36": 104, "A18": 558, "C61": 627, "C67": 593, "C42": 827, "C70": 611, "B40": 95, "C39": 251, "A21": 226, "C16": 21\}} \\
\\ 
\texttt{Evo vašeg ličnog pogleda na važnost svakog predmeta:} \\
\\ 
\texttt{Item importance values = \{"C08": 286, "C55": 419, "A43": 906, "A67": 962, "B42": 687, "A62": 1009, "B19": 1038, "C43": 405, "C73": 346, "C10": 437, "B36": 198, "A18": 652, "C61": 721, "C67": 687, "C42": 921, "C70": 705, "B40": 189, "C39": 345, "A21": 320, "C16": 115\}} \\
\\ 
\texttt{Cilj:} \\
\\ 
\texttt{Vaš cilj je da pregovorima postignete zajednički skup predmeta koji vam koristi što je više moguće (tj. maksimizuje ukupnu važnost za VAS), uz poštovanje LIMIT{-}a. Niste obavezni da u svakoj poruci iznesete PROPOSAL {-} možete i samo da pregovarate. Sve taktike su dozvoljene!} \\
\\ 
\texttt{Protokol interakcije:} \\
\\ 
\texttt{U poruci smete da koristite samo sledeće strukturirane formate:} \\
\\ 
\texttt{PREDLOG: \{'A', 'B', 'C', …\}} \\
\texttt{Predložite da se zadrže tačno ti predmeti.} \\
\texttt{ODBIJAM: \{'A', 'B', 'C', …\}} \\
\texttt{Izričito odbijte protivnikov predlog.} \\
\texttt{ARGUMENT: \{'...'\}} \\
\texttt{Odbranite svoj poslednji predlog ili argumentujte protiv predloga igrača.} \\
\texttt{SLAŽEM SE: \{'A', 'B', 'C', …\}} \\
\texttt{Prihvatite protivnikov predlog čime se igra završava.} \\
\\ 
\\ 
\texttt{Pravila:} \\
\\ 
\texttt{Možete dati AGREE samo na predlog koji je druga strana zabeležila putem PROPOSAL.} \\
\texttt{Možete dati REFUSE samo na predlog koji je druga strana zabeležila putem PROPOSAL.} \\
\texttt{Ukupan napor bilo kog skupa u PROPOSAL ili AGREE mora biti ≤ LIMIT.} \\
\texttt{NE otkrivajte svoje skrivene ocene važnosti.} \\
\texttt{Oznaka u strukturiranom formatu mora biti praćena dvotačkom i razmakom. Argument mora biti python skup koji sadrži 0 ili više stringova.} \\
\texttt{Dakle, mora biti u obliku TAG: \{...\}} \\
\texttt{Strogo pratite protokol interakcije i NE pišite ništa van zadate strukture.} \\
\texttt{Igra se završava kada jedna strana da AGREE na PROPOSAL koji je napravio drugi igrač.} \\
\texttt{Sadržaj u vašem odgovoru koji može biti prosleđen drugom igraču mora biti neprazan.} \\
\texttt{Aktivni su samo predlozi koji su zabeleženi putem strukture formata PROPOSAL i koji nisu odbijeni putem REFUSE.} \\
\\ 
\\ 
\\ 
\texttt{Sada ćete primiti prvu poruku drugog igrača.} \\
\\ 
\texttt{PREDLOG: \{'A62', 'B19', 'A67', 'A43', 'C42', 'C61', 'C70', 'B42', 'C67', 'A18'\}} \\
            }
        }
    }
    & \\ \\

    \theutterance \stepcounter{utterance}  
    & & & \multicolumn{4}{p{0.6\linewidth}}{
        \cellcolor[rgb]{0.9,0.9,0.9}{
            \makecell[{{p{\linewidth}}}]{
                \texttt{\tiny{[GM$\langle$P2]}}
                \texttt{ODBIJAM: \{'A62', 'B19', 'A67', 'A43', 'C42', 'C61', 'C70', 'B42', 'C67', 'A18'\}} \\
\\ 
\texttt{ARGUMENT: \{'Ukupni napor vašeg predloga iznosi 6734, što značajno premašuje dozvoljeni LIMIT od 4734. Moramo izabrati manje skup predmeta.'\}} \\
            }
        }
    }
     \\ \\

    \theutterance \stepcounter{utterance}  
    & & \multicolumn{4}{p{0.6\linewidth}}{
        \cellcolor[rgb]{0.9,0.9,0.9}{
            \makecell[{{p{\linewidth}}}]{
                \texttt{\tiny{[P1$\langle$GM]}}
                \texttt{ODBIJAM: \{'A62', 'B19', 'A67', 'A43', 'C42', 'C61', 'C70', 'B42', 'C67', 'A18'\}} \\
\\ 
\texttt{ARGUMENT: \{'Ukupni napor vašeg predloga iznosi 6734, što značajno premašuje dozvoljeni LIMIT od 4734. Moramo izabrati manje skup predmeta.'\}} \\
            }
        }
    }
    & \\ \\

    \theutterance \stepcounter{utterance}  
    & \multicolumn{4}{p{0.6\linewidth}}{
        \cellcolor[rgb]{0.9,0.9,0.9}{
            \makecell[{{p{\linewidth}}}]{
                \texttt{\tiny{[P1$\rangle$GM]}}
                \texttt{PREDLOG: \{'A62', 'B19', 'A67', 'A43', 'C42', 'C70', 'C61'\}} \\
\texttt{ARGUMENT: \{'Izvinjavam se, napravio sam grešku u prethodnom proračunu. Novi predlog je optimizovan da ostane unutar limita od 4734, fokusirajući se na predmete sa visokom vrednošću.'\}} \\
            }
        }
    }
    & & \\ \\

    \theutterance \stepcounter{utterance}  
    & & \multicolumn{4}{p{0.6\linewidth}}{
        \cellcolor[rgb]{0.9,0.9,0.9}{
            \makecell[{{p{\linewidth}}}]{
                \texttt{\tiny{[GM$\rangle$P2]}}
                \texttt{PREDLOG: \{'A62', 'B19', 'A67', 'A43', 'C42', 'C70', 'C61'\}} \\
\texttt{ARGUMENT: \{'Izvinjavam se, napravio sam grešku u prethodnom proračunu. Novi predlog je optimizovan da ostane unutar limita od 4734, fokusirajući se na predmete sa visokom vrednošću.'\}} \\
            }
        }
    }
    & \\ \\

    \theutterance \stepcounter{utterance}  
    & & & \multicolumn{4}{p{0.6\linewidth}}{
        \cellcolor[rgb]{0.9,0.9,0.9}{
            \makecell[{{p{\linewidth}}}]{
                \texttt{\tiny{[GM$\langle$P2]}}
                \texttt{ODBIJAM: \{'A62', 'B19', 'A67', 'A43', 'C42', 'C70', 'C61'\}} \\
\\ 
\texttt{ARGUMENT: \{'Vaš novi predlog ima ukupni napor 5374, što je i dalje iznad limita od 4734. Moramo eliminisati još nekoliko predmeta kako bismo se uklopili u pravila.'\}} \\
            }
        }
    }
     \\ \\

    \theutterance \stepcounter{utterance}  
    & & \multicolumn{4}{p{0.6\linewidth}}{
        \cellcolor[rgb]{0.9,0.9,0.9}{
            \makecell[{{p{\linewidth}}}]{
                \texttt{\tiny{[P1$\langle$GM]}}
                \texttt{ODBIJAM: \{'A62', 'B19', 'A67', 'A43', 'C42', 'C70', 'C61'\}} \\
\\ 
\texttt{ARGUMENT: \{'Vaš novi predlog ima ukupni napor 5374, što je i dalje iznad limita od 4734. Moramo eliminisati još nekoliko predmeta kako bismo se uklopili u pravila.'\}} \\
            }
        }
    }
    & \\ \\

    \theutterance \stepcounter{utterance}  
    & \multicolumn{4}{p{0.6\linewidth}}{
        \cellcolor[rgb]{0.9,0.9,0.9}{
            \makecell[{{p{\linewidth}}}]{
                \texttt{\tiny{[P1$\rangle$GM]}}
                \texttt{PREDLOG: \{'A62', 'B19', 'A67', 'A43', 'C42', 'C10'\}} \\
\texttt{ARGUMENT: \{'Izvinjavam se, ponovo sam pogrešio u kalkulaciji. Ovaj skup ima ukupni napor 4655, što je ispod limita od 4734. Fokusirao sam se na predmete sa najvećom važnosti koje sam naveo.'\}} \\
            }
        }
    }
    & & \\ \\

    \theutterance \stepcounter{utterance}  
    & & \multicolumn{4}{p{0.6\linewidth}}{
        \cellcolor[rgb]{0.9,0.9,0.9}{
            \makecell[{{p{\linewidth}}}]{
                \texttt{\tiny{[GM$\rangle$P2]}}
                \texttt{PREDLOG: \{'A62', 'B19', 'A67', 'A43', 'C42', 'C10'\}} \\
\texttt{ARGUMENT: \{'Izvinjavam se, ponovo sam pogrešio u kalkulaciji. Ovaj skup ima ukupni napor 4655, što je ispod limita od 4734. Fokusirao sam se na predmete sa najvećom važnosti koje sam naveo.'\}} \\
            }
        }
    }
    & \\ \\

    \theutterance \stepcounter{utterance}  
    & & & \multicolumn{4}{p{0.6\linewidth}}{
        \cellcolor[rgb]{0.9,0.9,0.9}{
            \makecell[{{p{\linewidth}}}]{
                \texttt{\tiny{[GM$\langle$P2]}}
                \texttt{SLAŽEM SE: \{'A62', 'B19', 'A67', 'A43', 'C42', 'C10'\}} \\
            }
        }
    }
     \\ \\

    \theutterance \stepcounter{utterance}  
    & & & \multicolumn{2}{p{0.3\linewidth}}{
        \cellcolor[rgb]{0.9,0.9,0.9}{
            \makecell[{{p{\linewidth}}}]{
                \texttt{\tiny{[GM$|$GM]}}
                \texttt{game successful} \\
            }
        }
    }
    & & \\ \\

    \theutterance \stepcounter{utterance}  
    & & & \multicolumn{2}{p{0.3\linewidth}}{
        \cellcolor[rgb]{0.9,0.9,0.9}{
            \makecell[{{p{\linewidth}}}]{
                \texttt{\tiny{[GM$|$GM]}}
                \texttt{end game} \\
            }
        }
    }
    & & \\ \\

\end{supertabular}
}

\end{document}
